\section{Fixed virtual classes and the divisor cosection}\label{sec:fixed}

We compute the ordinary fixed virtual classes furnished by
Proposition~\ref{prop:fixed-pot}.  Throughout this section, a symbol
\(\Div^{P}(X)\) denotes the scheme of effective divisors whose line bundle
lies in the specified component of \(\Pic(X)\); when that component is a
point, this is the complete linear system \(|P|\).
The universal divisor is denoted by \(\mathcal C\), and \(\pi\) always
denotes projection along the surface \(X\).
All derived Hom complexes below include this derived projection.
We retain the two settings of Section~\ref{sec:derived}; in particular,
\(\chi(\mathcal O_X)=1\) in both, and we put
\[
 m_P=\frac{P(P-K)}2=\chi(P)-1.
\]

The calculation has two different forms. When \(p_g=q=0\), each
Picard component is a point and the based Picard tangent vanishes.
The divisor factors are then complete linear systems, and the fixed
loci are smooth; we compute their virtual classes from their
obstruction bundles. On the auxiliary surface the divisor lines can
vary in a Picard curve. We retain this variation in the obstruction
theory, localize one divisor class to the empty or canonical divisor,
and only then restrict the other factor and the insertions to the
resulting fiber. The rank-one family decomposition and point-ideal
comparison below provide the common input to both calculations.

\subsection{Rank-one quotients in families}

The divisor-and-point description and its tautological obstruction
factors appear for a trivial rank-one target on a simply connected
surface in \cite[Section~5.2]{OP}. We record the factorization in
arbitrary families and the obstruction-theory comparisons needed here.

We write \(X^{[r]}\) for the Hilbert scheme of length-\(r\)
zero-dimensional subschemes of \(X\).

\begin{lemma}[Divisor and point factorization]\label{lem:rank-one-family}
Let \(R\) be a noetherian parameter scheme and let
\(L_R\twoheadrightarrow Q\) be a flat family of rank-zero quotients of a
line bundle on \(X\times R\).  Its kernel has a unique factorization
\[
 S'=L_R(-\mathcal C)\mathcal I_Z,
\]
where \(\mathcal C\) is a relative effective Cartier divisor and
\(Z\subset X\times R\) is finite and flat over \(R\).
Both factors commute with arbitrary base change.  Consequently, the
rank-one Quot functor with specified divisor component and point length
\(r\) is represented by
\begin{equation}\label{fix:rank-one-product}
 U=\Div^P(X)\times X^{[r]},
 \qquad S'=L(-\mathcal C)\mathcal I_Z.
\end{equation}
The assertion includes nonreduced parameter schemes.
\end{lemma}

\begin{proof}
Untwist by \(L_R\), and write \(I\subset\mathcal O_{X\times R}\) for the
kernel.  It is flat over \(R\), because the source and quotient are flat.
Every geometric fiber \(I_u\) is torsion-free of rank one on a smooth
surface and has projective dimension at most one.  Locally choose a
surjection from a vector bundle \(G\) onto \(I\).  Its kernel \(E\) is
flat over \(R\) and has locally free geometric fibers, so is locally
free by the local criterion of flatness.  Thus
\[
 0\longrightarrow E\longrightarrow G\longrightarrow I\longrightarrow0
\]
is a resolution by vector bundles, and remains a resolution after any
base change.

Set \(M=\det G\otimes(\det E)^{-1}=\det I\).  If \(E\) has rank \(e\),
wedging a section of \(G\) with the image of \(\bigwedge^eE\) gives the
maximal-minor map \(G\to M\); it kills \(E\) and induces \(I\to M\).
On the locus where \(I\) is invertible it is the canonical determinant
identification.  These local maps glue.  To justify this last assertion
over a nonreduced base, the fiber criterion for local freeness shows
that the closed noninvertible locus has finite fibers over \(R\);
properness makes this locus finite over \(R\).
Any closed support with finite fibers in the smooth relative
surface has grade at least two locally: choose two elements of its
ideal which form a regular sequence on the surface fiber and lift
them; successive applications of the local flatness criterion show
that they remain a regular sequence in the total space and that the
successive quotients are base-flat.  Powers of these elements
annihilate any coherent sheaf with that support; the regular-sequence
calculation therefore makes its sheaf \(\mathcal{H}om\) and
\(\mathcal Ext^1\) into a line bundle vanish.  In particular
a section of a line bundle supported there is zero.  Hence two local
minor maps which agree on the invertible locus agree everywhere.

On each geometric fiber the resulting map \(I_u\to M_u\) is the usual
inclusion of a rank-one torsion-free sheaf into its determinant line:
it is injective, and its cokernel has finite support.  Fiberwise
injectivity between base-flat modules now implies that \(I\to M\) is
injective with base-flat cokernel.  After twisting by \(M^{-1}\) it
is an ideal \(J\subset\mathcal O_{X\times R}\), and
\(\mathcal O/J\) is base-flat with finite fibers.  Properness of \(X\)
makes its support finite over \(R\).  It therefore defines a finite
flat subscheme \(Z\), with \(I=M\mathcal I_Z\).

The original inclusion \(I\hookrightarrow\mathcal O_{X\times R}\)
extends uniquely to a map \(M\to\mathcal O_{X\times R}\).
Indeed the grade-two observation, applied to the support of
\(M/I=M\mathcal O_Z\), gives
\[
 \mathcal{H}om(M\mathcal O_Z,\mathcal O_{X\times R})=
 \mathcal Ext^1(M\mathcal O_Z,\mathcal O_{X\times R})=0.
\]
Applying \(\mathcal{H}om(-,\mathcal O_{X\times R})\) to
\(0\to I\to M\to M\mathcal O_Z\to0\) proves the extension and uniqueness.
Its restriction to every surface fiber is nonzero, hence injective.
The same flatness criterion makes it an injection with base-flat
cokernel.  It is therefore a relative effective Cartier divisor
inclusion, \(M=\mathcal O(-\mathcal C)\).

The determinant resolutions, minor maps, and short exact sequences
all survive base change, the last because their quotients are
base-flat.  The regular sequences used for uniqueness also survive
base change, because their successive quotients are base-flat.
Thus uniqueness identifies the resulting factorization after every
base change, including a nonnoetherian one.
Conversely, multiplying a relative Cartier ideal by a finite flat
point ideal gives a flat quotient: its quotient is an extension of
\(\mathcal O_{\mathcal C}\) by \(\mathcal O_Z(-\mathcal C)\), both
base-flat.  This construction inverts the preceding one.  More
explicitly, any other line extending \(I\) with relatively finite
cokernel is uniquely identified with \(M\): the same grade-two
extension argument extends the identity on \(I\) in both directions,
and uniqueness makes the composites identities.  This proves the
functorial factorization, and hence the asserted isomorphism of
parameter schemes.
\end{proof}

Apply the Lemma to the two weight summands of a fixed quotient.  Write
\(D_1=T\), \(D_2=H\), and \(T+H=D\) in the Picard scheme.  If their
point lengths are \(n_1,n_2\), comparison of the second Chern characters
gives
\[
 n_1+n_2+L_1T+L_2H-\frac{T^2+H^2}{2}
 =n+L_1D-\frac{D^2}{2}.
\]
Since \(A=L_1-L_2\), this is
\begin{equation}\label{eq:length}
 n_1+n_2=n-(D-A)H+H^2.
\end{equation}
These labels form a finite set on the fixed locus under consideration:
the Picard components, Hilbert polynomials, and point lengths are
locally constant there, and the fixed locus is of finite type.

\subsection{The point-ideal complexes}

The fixed locus has been expressed in divisor and point factors.
We next identify the point contribution to its obstruction theory;
this will separate the Hilbert-scheme tangent from the excess bundle.

For a line bundle \(P\), write
\(P^{[r]}=\pi_*\mathcal O_Z(P)\) for its tautological bundle
of rank \(r\) on \(X^{[r]}\).
For the universal ideal \(\mathcal I\), the trace-free endomorphism
complex means the fiber of
\(\operatorname{tr}:\RHom(\mathcal I,\mathcal I)
\to\RGamma(\mathcal O_X)\).

\begin{lemma}[Trace splitting and the point obstruction]\label{fix:trace}
On \(X^{[r]}\), with universal ideal \(\mathcal I\), there is a
quasi-isomorphism
\begin{equation}\label{fix:hilbert-splitting}
 \RHom(\mathcal I,\mathcal I)
 \simeq \RGamma(\mathcal O_X)\oplus T_{X^{[r]}}[-1].
\end{equation}
The two maps to the right-hand side are the trace and the inverse of
the trace-free Kodaira--Spencer isomorphism.
On \(\Div^P(X)\times X^{[r]}\), define
\begin{equation}\label{fix:point-bundle}
 F=\bigl(\pi_*\mathcal O_Z(K-\mathcal C)\bigr)^\vee.
\end{equation}
It is a vector bundle of rank \(r\), and relative Serre duality gives
\begin{equation}\label{fix:point-duality}
 \RHom(\mathcal O_Z,\mathcal O(\mathcal C))=F[-2].
\end{equation}
\end{lemma}

\begin{proof}
The scalar map \(\RGamma(\mathcal O_X)\to
\RHom(\mathcal I,\mathcal I)\) is split by trace, since \(\mathcal I\)
has rank one.  At a geometric point \(Z\), extending across its
codimension-two support gives
\[
 \Hom(I_Z,I_Z)=\mathbb C,\qquad
 \Hom(I_Z,I_Z(K))=H^0(K).
\]
For example, an endomorphism extends to the determinant line and is
there multiplication by a global section; conversely multiplication
preserves the ideal.  By Serre duality the trace-free
\(\Ext^2(I_Z,I_Z)\) vanishes, as does the trace-free \(\Ext^0\).
Riemann--Roch gives
\(\chi(I_Z,I_Z)=\chi(\mathcal O_X)-2r\); hence the trace-free
\(\Ext^1\) has dimension \(2r\).

The boundary map of the point-ideal sequence embeds
\[
 T_ZX^{[r]}=\Hom(I_Z,\mathcal O_Z)
       \longrightarrow \Ext^1(I_Z,I_Z).
\]
Indeed every map \(I_Z\to\mathcal O_X\) is a scalar multiple of
inclusion, so its composite with \(\mathcal O_X\to\mathcal O_Z\)
is zero.  This boundary is the Kodaira--Spencer map of the universal
ideal.  Its trace vanishes, because the universal point ideal has
constant determinant.  Fogarty's smoothness Theorem
\cite[Theorem~2.4]{Fogarty} gives \(\dim X^{[r]}=2r\), so this is an isomorphism
onto the trace-free part.  The family complexes are perfect and
commute with derived base change.  A map between them is a
quasi-isomorphism if it is so on geometric fibers, proving
\eqref{fix:hilbert-splitting}.

For the other assertion, \(Z\) is finite flat of length \(r\), so
\(\pi_*\mathcal O_Z(K-\mathcal C)\) is locally free of that rank,
with arbitrary base change and no higher direct images.  Relative
Serre duality for the smooth proper surface projection identifies
the dual of \(\RHom(\mathcal O_Z,\mathcal O(\mathcal C))\) with
\(\pi_*\mathcal O_Z(K-\mathcal C)[2]\).
Dualizing gives precisely \eqref{fix:point-duality}, including its
shift and the indicated dual bundle.
\end{proof}

For later calculations it follows that the virtual tangent class of
a rank-one arrow is
\begin{equation}\label{fix:arrow-class}
 [T_U^{\vir}]
 =[R\pi_*\mathcal O(\mathcal C)]-[F]
       -[\RGamma(\mathcal O_X)]+[T_{X^{[r]}}].
\end{equation}
In fact it is the cone of
\(\RHom(\mathcal I,\mathcal I)\to
\RHom(\mathcal I,\mathcal O(\mathcal C))\), and the ideal sequence
and \eqref{fix:point-duality} give the formula.

\begin{lemma}[Euler class for a smooth morphism]\label{fix:smooth-euler}
Suppose \(g:Y\to Z\) is smooth and carries a relative perfect
obstruction theory \(E_g\to L_{Y/Z}\) of amplitude \([-1,0]\).
Then \(\Ob_g=h^1(E_g^\vee)\) is a vector bundle,
\[
 [\Ob_g]=[T_{Y/Z}]-[E_g^\vee]\quad\hbox{in }K^0(Y),
\]
and its virtual pullback is
\begin{equation}\label{fix:smooth-operation}
 g^!\gamma=\eu(\Ob_g)\cap g^*\gamma.
\end{equation}
The formula holds over arbitrary base cycles.  For \(Z\) a point,
it says \([Y]^\vir=\eu(\Ob_Y)\cap[Y]\).
\end{lemma}

\begin{proof}
Here \(L_{Y/Z}=\Omega_{Y/Z}\) is a vector bundle in degree zero.
Locally split the surjection from a two-term presentation of \(E_g\)
onto this bundle.  The remaining acyclic degree-zero terms can be
cancelled; one obtains
\[
 E_g\simeq \Omega_{Y/Z}\oplus\Ob_g^\vee[1],
\]
with comparison map the projection.  This shows local freeness
and the stated \(K\)-class, without requiring a global splitting.
The relative intrinsic normal cone of a smooth map is
\([0/T_{Y/Z}]\).  The dual comparison morphism gives its global
inclusion into \(h^1/h^0(E_g^\vee)\), with quotient bundle
\(\Ob_g\).  Locally this is the inclusion into
\([\Ob_g/T_{Y/Z}]\) supplied by the splitting above; globally the
subbundle excess formula gives its zero-section Gysin image as
the Euler class of \(\Ob_g\).  Applying this excess formula over
each base cycle proves \eqref{fix:smooth-operation}.
\end{proof}

\subsection{The fixed class when \texorpdfstring{\(p_g=q=0\)}{p\_g=q=0}}

\begin{proposition}\label{prop:fixed-pg0}
Assume \(p_g=q=0\).  A nonempty fixed part with lines \(T,H\)
and lengths satisfying \eqref{eq:length} is the smooth product
\[
 W=|T|\times|H|\times X^{[n_1]}\times X^{[n_2]}.
\]
Let \(x_i=c_1(\mathcal O_{|D_i|}(1))\), with \(D_1=T,D_2=H\).
The fixed virtual class is
\begin{equation}\label{fix:pg0-class}
 [W]^\vir=
 \prod_{i=1}^2\bigl(x_i^{h^1(D_i)}\eu(F_i)\bigr)\cap[W],
 \qquad
 F_i=\bigl((K-D_i)^{[n_i]}\bigr)^\vee\otimes\mathcal O(1_i).
\end{equation}
A part with either \(m_T<0\) or \(m_H<0\) contributes zero.
For the other parts, after including \(\eu(F_1)\eu(F_2)\),
pushforward over the linear systems is coefficient extraction at
\([x_1^{m_T}x_2^{m_H}]\).
\end{proposition}

\begin{proof}
The Picard scheme is discrete and reduced, and
\(\RGamma(\mathcal I_o)=0\).  Thus the determinant condition
leaves just the indicated products, with the sum of their two
rank-one obstruction theories.  For an effective line \(P\),
multiplication by any nonzero section of \(P\) injects
\(H^0(K-P)\) into \(H^0(K)=0\), so \(H^2(P)=0\).
On its linear system,
\(\mathcal O(\mathcal C)=P\boxtimes\mathcal O(1)\).
The Euler sequence and \eqref{fix:arrow-class} therefore give
\[
 [T_{|P|\times X^{[r]}}]-[T_U^\vir]
 = [H^1(P)\otimes\mathcal O(1)]+[F].
\]
Apply Lemma~\ref{fix:smooth-euler}.  Equality of these vector
bundle classes gives equality of their Chern classes, and hence
\eqref{fix:pg0-class}; it does not assert a global splitting of
the obstruction bundle.

Finally,
\[
 \dim|P|-h^1(P)=h^0(P)-1-h^1(P)=m_P.
\]
If \(m_P<0\), the factor \(x^{h^1(P)}\) is already zero in
\(\CH^*(|P|)\).  Otherwise integration against that factor
replaces extraction at degree \(\dim|P|\) by extraction at degree
\(m_P\).  The localized moving factors are expanded with nonzero
equivariant constant terms, so introduce no negative hyperplane
powers.  Thus this argument applies also to the localized
insertions used below.
\end{proof}

\subsection{A compatible relative theory in the auxiliary setting}

We now assume \(p_g=q=1\), that \(K\) is ample, and that \(D=3K\).
We impose the determinant constraint in one rank-one factor and apply
the divisor cosection to the other.
On each fixed part choose an index \(a\) such that
\(KD_a>K^2\), which is possible since
\(KD_1+KD_2=3K^2\).  Call it the high index, and call the other
index \(b\) the low index.  If both satisfy the inequality,
choose either one.  The ensuing vanishing argument applies to
the resulting low index as well.

\begin{lemma}[Relative determinant pullback]\label{fix:relative-pot}
For the projection \(f:W\to U_b\) to the unrestricted low
rank-one Quot, there is a relative perfect obstruction theory
whose virtual tangent is
\begin{equation}\label{fix:relative-tangent}
 T_f^\vir=
 \fib\!\left\{\RHom(S'_a,Q_a)
       \longrightarrow\RGamma(\mathcal I_o)[1]\right\}.
\end{equation}
It is compatible with the absolute theories of \(W\) and \(U_b\),
and
\begin{equation}\label{fix:relative-class}
 [W]^\vir=f^![U_b]^\vir.
\end{equation}
\end{lemma}

\begin{proof}
Use the actual projection of enhancements
\(\mathfrak W\to\mathfrak U_b\) from
Proposition~\ref{prop:fixed-pot}.  The tangent of its determinant
fiber in the high direction is \eqref{fix:relative-tangent}.
The map there is the high arrow's boundary followed by trace
and based normalization.  Lemma~\ref{der:trace-surjectivity},
with rank one and \(KD_a>K^2\), makes its degree-one map onto
\(H^2(\mathcal O_X)\) surjective, so the fiber has amplitude
\([0,1]\).  Explicitly, an annihilating canonical form would
annihilate \(Q_a\); at each generic curve point a quotient of
one generator annihilated by that form has length at most its
vanishing order.  Thus \(KD_a\le K^2\), a contradiction.

Restrict the derived cotangent transitivity triangle to \(W\).
It has a natural morphism to the classical transitivity
triangle, giving
\[
\begin{array}{ccccc}
 Lf^*E_{U_b}&\longrightarrow&E_W&\longrightarrow&E_f\\
 \downarrow&&\downarrow&&\downarrow\\
 Lf^*L_{U_b}&\longrightarrow&L_W&\longrightarrow&L_{W/U_b}.
\end{array}
\]
Each row continues by a shift of its first term.
Here the absolute comparisons have cones in degrees at most
\(-2\), by Lemma~\ref{der:truncation}.  If these cones are
\(C_{U_b}\) and \(C_W\), and the relative comparison cone is
\(C_f\), the diagram gives a triangle
\[
 Lf^*C_{U_b}\longrightarrow C_W\longrightarrow C_f
       \longrightarrow Lf^*C_{U_b}[1].
\]
Derived pullback preserves the upper cohomological bound, and
the shift \([1]\) lowers that bound.  Hence \(C_f\) is also
concentrated in degrees at most \(-2\).  Combined with the
amplitude calculation, this proves that \(E_f\) is a relative
perfect obstruction theory.

We have thus produced the compatible triple, including its
morphisms to the cotangent triangle, required in
\cite[Definition~4.5]{Manolache}.  The morphisms are between
quasiprojective schemes, so are of Deligne--Mumford type and
have the required stratifications by global quotients.
Virtual-pullback functoriality
\cite[Theorem~4.8 and Corollary~4.9]{Manolache} gives
\eqref{fix:relative-class}.
\end{proof}

\subsection{The point obstruction as an excess bundle}

Let \(\mathfrak D^P\) denote the divisor enhancement: the open
rank-one arrow stack with invertible source and divisor
quotient.  Its classical truncation is \(\Div^P(X)\).
For a divisor \(C\), with source \(M=L(-C)\), its tangent is
\[
 \RHom(M,L/M)=R\pi_*\mathcal O_C(C).
\]
It has amplitude \([0,1]\), and its Euler characteristic is
\(\chi(P)-\chi(\mathcal O_X)=m_P\).  Thus it induces a divisor
virtual class of dimension \(m_P\).

\begin{lemma}[Compatible point excess]\label{fix:excess}
On the product \(U_b=\Div^{D_b}(X)\times X^{[n_b]}\), the
rank-one Quot virtual class is
\begin{equation}\label{fix:low-product}
 [U_b]^\vir=
 \eu(F_b)\cap
       \bigl([\Div^{D_b}(X)]^\vir\times[X^{[n_b]}]\bigr).
\end{equation}
\end{lemma}

\begin{proof}
There is a morphism of derived moduli problems
\begin{equation}\label{fix:product-enhancement-map}
 g:\mathfrak D^{D_b}\times X^{[n_b]}\longrightarrow\mathfrak U_b,
 \qquad
 (M_b\to L_b,\mathcal I_b)\longmapsto
       (M_b\mathcal I_b\to M_b\to L_b).
\end{equation}
The ordinary smooth Hilbert scheme here carries its smooth
enhancement, and its universal ideal is perfect.
Tensoring and composing perfect families define the displayed
map on derived bases as well.  Lemma~\ref{lem:rank-one-family}
identifies its truncation with the identity of \(U_b\).

To compute its differential, first forget the arrows and
compare their source moduli.  On the common truncation the map is
\begin{equation}\label{fix:source-tangent-map}
 \RGamma(\mathcal O_X)[1]\oplus T_{X^{[n_b]}}
       \longrightarrow
       \RHom(\mathcal I_b,\mathcal I_b)[1].
\end{equation}
The first summand acts by scalars and the second is the
Kodaira--Spencer map of the ideal, so
Lemma~\ref{fix:trace} makes this a quasi-isomorphism.
Indeed, on each geometric fiber the scalar map identifies the trace
summands, and the Hilbert map induces an isomorphism onto the
trace-free \(\Ext^1\).  The source derivative is therefore triangular
with isomorphic diagonal blocks.  Its perfect cone has zero geometric
fibers, so vanishes; no choice of a global splitting is needed here.
The scalar assertion follows directly by differentiating
tensoring with a line.  Before taking sections it is
\(\mathcal O_X[1]\to
R\mathcal{H}om(\mathcal I_b,\mathcal I_b)[1]\),
and corresponds to the identity endomorphism.

The relative arrow tangent map over these source moduli is
\[
 \RHom(M_b,L_b)\longrightarrow
       \RHom(M_b\mathcal I_b,L_b).
\]
The point-ideal sequence and
\eqref{fix:point-duality} identify its cone with
\[
 \RHom(\mathcal O_{Z_b},\mathcal O(\mathcal C_b))[1]
       =F_b[-1].
\]
The comparison of the two source tangent triangles, using
\eqref{fix:source-tangent-map}, therefore proves
\begin{equation}\label{fix:excess-triangle}
 T_{\mathrm{prod}}^\vir\longrightarrow T_{U_b}^\vir
       \longrightarrow F_b[-1]
       \longrightarrow T_{\mathrm{prod}}^\vir[1].
\end{equation}

Crucially, dualizing this is a morphism of perfect obstruction
theories \(E_{U_b}\to E_{\mathrm{prod}}\), compatible with
their maps to the same \(L_{U_b}\), because it came from
\eqref{fix:product-enhancement-map}.  Thus the induced map
of obstruction bundle stacks is an inclusion
\[
 h^1/h^0(T_{\mathrm{prod}}^\vir)
       \longhookrightarrowactual
 h^1/h^0(T_{U_b}^\vir)
\]
with quotient vector bundle \(F_b\).  For precision, on an
affine open the extension in \eqref{fix:excess-triangle}
splits: its class lies in
\(\Hom(F_b,T_{\mathrm{prod}}^\vir[2])=0\), since \(F_b\)
is locally free and the tangent has amplitude \([0,1]\).
This describes the inclusion locally as a vector bundle
substack; the quotient is globally the displayed \(F_b\).

Compatibility of the obstruction morphisms implies that the
intrinsic normal cone embedded using the Quot theory is the
image of its embedding using the product theory.
The subbundle excess formula for zero-section Gysin maps
therefore multiplies the latter virtual class by \(\eu(F_b)\).
One can see the formula on any base cycle by pulling to the
larger bundle stack: the smaller stack is the zero locus of
the quotient-bundle coordinate, so its inclusion multiplies
the pullback class by the top Chern class of that quotient.
This is the zero-section excess formula of
\cite[\S6.3]{Fulton}; its vector-bundle form is also
\cite[Lemma~2.5]{KiemLi}.  We apply it to the bundle-stack presentations
of the intrinsic cone in \cite{BF}.
The product theory has class
\([\Div^{D_b}(X)]^\vir\times[X^{[n_b]}]\).
This proves \eqref{fix:low-product} on the possibly singular
divisor scheme.
\end{proof}

\subsection{Localization on the divisor scheme}

The excess comparison separates the low point factor from its divisor
virtual class. We now determine the support of that divisor class
and the sign of its canonical contribution, using the divisor cosection
studied by Chang--Kiem \cite[Section~3.2]{ChangKiem}.

\begin{lemma}[Support of the divisor class]\label{fix:cosection-support}
In the auxiliary setting a divisor component with \(m_P<0\)
has zero virtual class.  If \(m_P\ge0\), its virtual class
is supported, by cosection localization, on the empty
divisor or on the unique canonical divisor \(z\).
The empty-divisor class is the reduced point with coefficient
one.  Only the Picard component containing \(K\) can give
the canonical contribution.
\end{lemma}

\begin{proof}
The first assertion is the vanishing of a Chow group in
negative dimension.  For the second, the universal divisor
sequence yields a cosection
\begin{equation}\label{fix:cosection}
 \sigma:\Ob_{\Div^P}
 =R^1\pi_*\mathcal O_{\mathcal C}(\mathcal C)
       \longrightarrow H^2(\mathcal O_X)\otimes\mathcal O
       \simeq\mathcal O.
\end{equation}
At a divisor \(C\) of line class \(P\) the long exact
cohomology sequence ends in
\[
 H^1(\mathcal O_C(C))\xrightarrow{\sigma_C}H^2(\mathcal O_X)
       \longrightarrow H^2(P)\longrightarrow0.
\]
The obstruction fibers agree with these groups because the
perfect tangent has amplitude \([0,1]\).  Thus
\(\coker(\sigma_C)=H^2(P)=H^0(K-P)^\vee\).
The cosection is not surjective exactly where \(K-P\) is
effective.  Kiem--Li localization
\cite[Theorem~1.1]{KiemLi} expresses the ordinary virtual
class as the pushforward of a class on this degeneracy locus.

Suppose a degeneracy point exists and put \(r=KP\).
Both \(P\) and \(K-P\) are effective.  Since \(K\) is ample
and \(m_P\ge0\), Hodge index gives
\[
 0\le r\le K^2,\qquad
 r\le P^2\le\frac{r^2}{K^2}.
\]
For \(0<r<K^2\) these inequalities contradict one another.
If \(r=0\), ampleness forces the effective divisor \(C\)
to be empty.  If \(r=K^2\), any effective divisor in
\(K-P\) has zero intersection with \(K\), so is empty.
In the latter case \(P\cong K\) as line bundles, and
\(h^0(K)=1\) gives exactly one divisor \(z\).

For the empty divisor the tangent complex
\(R\pi_*\mathcal O_\varnothing\) is zero.
Its parameter scheme is therefore a reduced point with
the ordinary point class.  A nontrivial numerically
trivial line has no effective divisor, again by ampleness.
This proves all the support assertions.
\end{proof}

\begin{lemma}[Canonical multiplicity and sign]\label{fix:canonical-sign}
For the divisor component whose Picard component contains
\(K\), there is an integer \(s>0\) such that
\begin{equation}\label{eq:canonical-cycle}
 [\Div^{K}(X)]^\vir=-s[z].
\end{equation}
Here \(\Div^{K}(X)\) allows variation of the line in its
Picard component; \(z\) is the unique effective canonical
divisor.
\end{lemma}

\begin{proof}
Let \(C_{\Pic}\) be the Picard component containing \(K\);
it is a smooth projective curve because \(q=1\).
Choose an open neighborhood \(V\) of the point \(K\).
Upper semicontinuity, followed by shrinking, gives
\(h^0(P)\le1\) for \(P\in V\).
For \(P\ne K\) in this component, \(K-P\) is a nontrivial
numerically trivial line, so \(h^2(P)=0\).
Riemann--Roch gives \(\chi(P)=1\), and consequently
\(h^0(P)=1+h^1(P)\ge1\).
At \(K\), \(h^0(K)=1\) by assumption.
Thus \(h^0(P)\) is constantly one on \(V\).

We verify the resulting scheme structure, including
infinitesimal families.  Represent the derived direct image
of the based universal line over \(V\), locally, by vector
bundles
\[
 E^0\xrightarrow{d^0}E^1\xrightarrow{d^1}E^2.
\]
There are no negative terms because this is the direct image
of a line bundle, and relative dimension two bounds the
positive terms.  The kernel of the fiber of \(d^0\) has
constant dimension one.  Since \(V\) is reduced, all minors
above the constant rank vanish scheme-theoretically.
After inverting a maximal nonzero minor, elementary row
and column operations turn \(d^0\) into an identity block
and a zero block.  Its kernel is therefore a line
subbundle and commutes with arbitrary, including
nonreduced, base change.  It is the bundle of sections of
the universal line.  The divisor functor over \(V\) is
its projective bundle of nonzero sections up to scalar,
and hence is \(V\) itself.  Thus the divisor scheme is
a smooth curve in a neighborhood of \(z\).

The divisor theory has virtual dimension \(m_K=0\).
Lemma~\ref{fix:smooth-euler} makes its obstruction sheaf
an obstruction line \(\mathcal E\) on this neighborhood.
At \(z\) the map \(H^2(\mathcal O_X)\to H^2(K)\) in
the divisor sequence is an isomorphism: its dual maps
\(1\) to the nonzero canonical section.
Hence the cosection \(\mathcal E\to\mathcal O_V\)
vanishes at \(z\).  By Lemma~\ref{fix:cosection-support}
it is nonzero at every other point nearby.  On the smooth
curve its germ is multiplication by a unit times a
uniformizer to a finite positive power \(s\).
It follows that the line with its cosection is isomorphic to
\[
 \bigl(\mathcal O_V(-sz),\
       \mathcal O_V(-sz)\hookrightarrow\mathcal O_V\bigr).
\]

Shrink to an affine neighborhood of \(z\).
As in Lemma~\ref{fix:smooth-euler}, the perfect
obstruction theory there splits, with its comparison map, as
\(\Omega_V\oplus\mathcal E^\vee[1]\to\Omega_V\).
The localized intrinsic-cone calculation is thus the
localized Euler class of the zero section in
\(\mathcal E=\mathcal O_V(-sz)\).
For clarity, in the regularizing formula
\cite[Definition~2.3 and equation~(2.1)]{KiemLi},
take the identity modification of \(V\), the quotient line
\(\mathcal O_V(D)\) with \(D=-sz\), and its zero kernel.
The lifted zero cone is \(V\), so the formula gives
\[
 [D]\cap[V]=-s[z].
\]
This determines the coefficient of the global localized class:
localized Gysin maps commute with restriction to an open
neighborhood of the degeneracy locus, by restricting the
regularizing modification, its cone cycles, and the ordinary
Gysin maps. This excision is also stated in
\cite[Proposition~2.3]{ChangKiem}.
The support lemma shows that \(z\) is the only degeneracy
point on the whole divisor component. Pushing its localized
class forward, by \cite[Theorem~1.1]{KiemLi}, proves
\eqref{eq:canonical-cycle}.
The argument used smoothness of the Picard curve; it imposed no
smoothness or reducedness condition on the canonical divisor itself.
\end{proof}

Chang and Kiem determine the canonical-divisor virtual coefficient more
precisely as $(-1)^{\chi(\mathcal O_X)}$ for minimal surfaces of general
type with $p_g>0$ \cite[Theorem~1.3 and Corollary~3.4]{ChangKiem}.
Their divisor theory has obstruction sheaf
\(R^1\pi_*\mathcal O_{\mathcal C}(\mathcal C)\)
\cite[Theorem~3.1]{ChangKiem}, and its cosection is
\eqref{fix:cosection} after identifying \(H^2(\mathcal O_X)\) with
\(\mathbb C\) using a nonzero canonical form
\cite[Section~3.2]{ChangKiem}.
On the smooth neighborhood \(V\) above, both theories therefore give
the localized Euler class of the same obstruction line and cosection.
Excision identifies their coefficients, so \(s=1\) here.
The local argument above retains only the sign and nonvanishing needed
here.  This sign comes from the divisor cosection; the moving Euler
classes have their own weight signs, computed separately in
Section~\ref{sec:relation}.

\subsection{The high factor and restriction to its fibers}

\begin{proposition}\label{prop:fixed-auxiliary}
Assume \(p_g=q=1\), \(K\) ample, and \(D=3K\).
Every fixed contribution vanishes except those obtained
from a low divisor equal to the empty divisor or to the
canonical divisor.  In either surviving case the
contribution is computed on the actual fiber
\[
 |D_a|\times X^{[n_a]}\times X^{[n_b]},
\]
with the factors \(\eu(F_a)\eu(F_b)\).
The low hyperplane is \(x_b=0\) and \(m_{D_b}=0\);
pushforward over \(|D_a|\) is extraction at
\(x_a^{m_{D_a}}\).
The overall multiplier is \(1\) for the empty divisor and
\(-s\), with the same \(s>0\) as in
\eqref{eq:canonical-cycle}, for the canonical divisor.
The statement applies to operational insertions in
universal perfect complexes, also after adjoining marked
surface factors.
\end{proposition}

\begin{proof}
By Lemmas~\ref{fix:relative-pot} and \ref{fix:excess},
the fixed class is obtained by applying \(f^!\) to
\[
 \eu(F_b)\cap
 \bigl([\Div^{D_b}(X)]^\vir\times[X^{[n_b]}]\bigr).
\]
The support and canonical-sign Lemmas make this input
zero unless its divisor class is the point class of the
empty divisor or \(-s[z]\).

Over each of the corresponding low Picard components,
the high line determined by the determinant has
\(D_a-K\) numerically \(2K\) or \(K\), respectively.
It is therefore ample.  Kodaira vanishing gives
\(H^j(D_a)=0\) for \(j>0\), uniformly in that component.
Use the based universal low line to obtain the based high
line as \(D-D_b\).  Its pushforward is a vector bundle
of rank \(\chi(D_a)=1+m_{D_a}\), with arbitrary base
change.  Lemma~\ref{lem:rank-one-family} consequently
expresses \(f\) here as its relative projective bundle
of high divisors, followed by the product with the smooth
Hilbert scheme \(X^{[n_a]}\).  In particular \(f\) is smooth,
even if the low divisor scheme away from the degeneracy
point is singular.

Formula \eqref{fix:relative-tangent}, together with
\eqref{fix:arrow-class}, gives
\[
\begin{aligned}
 [T_f^\vir]
  ={}&[R\pi_*\mathcal O(\mathcal C_a)]-[F_a]
        +[T_{X^{[n_a]}}]\\
    &-[\RGamma(\mathcal O_X)]
        -[\RGamma(\mathcal I_o)[1]].
\end{aligned}
\]
The evaluation sequence at \(o\) implies
\[
 [\RGamma(\mathcal O_X)]
       +[\RGamma(\mathcal I_o)[1]]=1.
\]
Thus the two determinant and scalar terms together
subtract exactly one scalar.  Higher cohomology vanishes
for the high line, and the relative projective Euler
sequence now identifies
\[
 [T_f^\vir]=[T_f]-[F_a].
\]
Lemma~\ref{fix:smooth-euler} proves, on all base cycles,
\begin{equation}\label{fix:high-operation}
 f^!\gamma=\eu(F_a)\cap f^*\gamma.
\end{equation}
Here the use of the \(K\)-class is legitimate: the
compatible relative perfect obstruction theory has
already been constructed, and smoothness identifies
its virtual operation with the Euler class of a vector
bundle.  Its Chern classes are determined by the
displayed \(K\)-class.

Apply \eqref{fix:high-operation} to the point-supported
low input class.  Compatibility of smooth pullback with
closed pushforward restricts it to the fiber over that
low divisor point, while retaining \(X^{[n_b]}\).
The projection formula similarly restricts every
operational insertion to the fiber.  The universal
divisors there are \(D_i\boxtimes\mathcal O(1_i)\),
with the low factor a point; in particular \(x_b=0\).
The high projective space has dimension
\(h^0(D_a)-1=m_{D_a}\), proving the extraction rule.
The universal point ideals, quotient complexes, and
their derived Hom complexes restrict by derived base
change for the smooth proper surface projection.
All of these operations commute with adjoining the
marked surface factors.  This proves the asserted
fiber formula with the multipliers already calculated.
\end{proof}
